\section{Coverage of the source articles}\label{sec:coverage}
The table follows MathWorld's equation numbers \cite{mw}. The numbered
claims are all covered. The entries on the right identify the arithmetic
content of the corresponding infinite formulas.

\begin{center}\small
\begin{tabular}{@{}p{.13\linewidth}p{.40\linewidth}p{.40\linewidth}@{}}
\toprule
Equations & Statement & Proof or interpretation\\\midrule
(1), (2) & Definition of $\mu$, summatory $M$ &
Valuations and bounded sums, \S\ref{sec:core}.\\
(3), (4) & Reciprocal Dirichlet series and Euler factors &
Positive half-plane margin, prime-mask expansion, \S\ref{sec:dirichlet-interior}.\\
(5), (12) & Lambert series and the exponential product &
Truncated coefficients and analytic error bounds, \S\ref{sec:mob-lambert}.\\
(6), (7) & The first two moments at $s=1$ &
Ordinary partial sums with a PR modulus, \S\ref{sec:boundary}.\\
(8)--(10) & The three reciprocal-square masses &
Finite trigonometric proof of the two even zeta values, \S\ref{sec:mob-basel}.\\
(11), (13) & Squarefree divisor count and product law &
Prime masks and additive valuations, \S\ref{sec:core}.\\
(14), (15) & Divisor cancellation and $\mu*\tau=\mathbf1$ &
The involution and convolution associativity, \S\ref{sec:core}.\\
\bottomrule
\end{tabular}
\end{center}

The additional formulas on Wikipedia \cite{wiki} include the Liouville
description, $\zeta(s)/\zeta(2s)$, the prime-shifted Lambert series,
the second logarithmic moment at $s=1$, the trigonometric unital sum
\cite{kline}, the two primitive-root identities, the Farey sum, the
irreducible-polynomial count, and convolution powers with a polynomial
parameter. Sections~\ref{sec:core}--\ref{sec:pnt} prove these and the
mean-value equivalence with PNT. Section~\ref{sec:boundary} also proves a conditional zero-exclusion bound from power cancellation of $M$. A computable finite-interval presentation
is specified for the incidence-algebra generalization. The occupation
number argument explains the arithmetic content of the gas model.

The value lists, including all sixty tabulated values, the initial sign
fibers, and the displayed products of three and five distinct primes,
are reproduced by \texttt{checks.py}. For example, bounding the scan
by 9690 certifies completeness through the last listed product of five
primes. The script also checks independent sample counts over two finite
fields. The published plots are illustrations of the same decidable
tests. The articles' historical notes, software listings, and references
to other problems have no additional theorem statement about $\mu$ to
prove here. Comparing the complexity of the two primitive-root and
prime-factor evaluation procedures would require a specified computational
model and representation of the output. We give exact bounded algorithms
for both and make no complexity comparison.

\begin{thebibliography}{9}\raggedright
\bibitem{wiki} Wikipedia contributors, \emph{Möbius function}, article
used for the formula inventory. \href{https://en.wikipedia.org/wiki/M\%C3\%B6bius_function}{Article}.
\bibitem{mw} E. W. Weisstein, \emph{Möbius Function}, MathWorld,
equations (1)--(15). \href{https://mathworld.wolfram.com/MoebiusFunction.html}{Article}.
\bibitem{kline} J. Kline, \emph{Unital Sums of the Möbius and Mertens Functions},
Journal of Integer Sequences 23 (2020), article 20.8.1.
\href{https://cs.uwaterloo.ca/journals/JIS/VOL23/Kline/kline4.html}{Paper}.
\bibitem{proveit} V. Reshetnikov et al., \emph{ProveIt},
\texttt{Logic/PrimitiveRecursiveArithmetic}, in particular
\texttt{ConciseIPRA} and \texttt{Derivations}. Declarations inspected at
\texttt{65b76a5b8aa8} (full revision in the link).
\href{https://github.com/VladimirReshetnikov/ProveIt/tree/65b76a5b8aa87ecad4103c2d287eff5660d3c84c/Logic/PrimitiveRecursiveArithmetic}{Pinned directory}.
\bibitem{strip} ProveIt, \emph{Logarithmic-integral strengthenings of the
prime number theorem in iPRA}, November 15, 2026.
The files \texttt{tex/analytic.tex} (E1)--(E8), (A1)--(A7),
\texttt{tex/finite\_contour.tex}, and \texttt{tex/theta\_bridge.tex}
supply the quantitative analytic lemmas reused in
Sections~\ref{sec:boundary} and~\ref{sec:bootstrap}.
These are human proof sources. The new signed smoothing and the boundary
moment deductions are given in this report.
\href{https://github.com/VladimirReshetnikov/ProveIt/tree/a9c0f1949be61679ed2e219b97dc42ecb7ae91a9/Logic/PrimitiveRecursiveArithmetic/docs/proofs/analytic-number-theory/pnt_li_ipra_2026-11-15}{Pinned packet}.
\end{thebibliography}
